% introduction.tex

\section{INTRODUCTION} \label{sec:introduction}

This laboratory investigates a spectrum of anomaly detection techniques --- from shallow one-class classifiers to deep autoencoders and purely unsupervised clustering --- applied to a network traffic dataset derived from the KDD Cup 1999 benchmark.
The dataset contains 43 features describing individual network connections, with labels indicating the traffic type (Normal, DoS, Probe, R2L).
Our analysis is structured around four main tasks:
\begin{enumerate}
    \item \textbf{Dataset characterization and preprocessing} (\Cref{sec:task1}): We explore the dataset structure, handle categorical and numerical features, and analyze per-class feature distributions through heatmaps.
    \item \textbf{Shallow anomaly detection} (\Cref{sec:task2}): We evaluate One-Class SVM under three training configurations: normal data only (novelty detection), all data (outlier detection), and mixed contamination levels. We quantify the impact of the $\nu$ parameter and training data composition.
    \item \textbf{Deep anomaly detection and data representation} (\Cref{sec:task3}): We train an autoencoder exclusively on normal data and explore three detection strategies: reconstruction error thresholding, OC-SVM on bottleneck embeddings, and OC-SVM on PCA-reduced features.
    \item \textbf{Unsupervised anomaly detection and interpretation} (\Cref{sec:task4}): We apply K-Means and DBSCAN to assess whether purely unsupervised approaches can identify attack patterns without any label supervision, using silhouette analysis, t-SNE visualization, and cluster purity metrics.
\end{enumerate}
